Most of the present chapter recaps known results.
However, the results on one-point singulars,
for example, \cref{prop:char-one-attractor,prelim:thm::one-attractor-idempotent,prelim:prop:one-attractor-idempotent-fin},
and those on compositions of 
our classes of transformations, are, to the best of our knowledge, new.
One-point singulars were introduced
in~\cite{Rystsov2024}, but this type of transformations
has been investigated before,
for example in Ara\'ujo \& Cameron~\cite{AraujoCameron2014} on finite sets  where
they are referred to via their kernel type $(k,1,\ldots,1)$ with $k \ge 2$.
We are also not aware of any studies that investigated
when the composition of two transformations results in a one-point singular,
or an idempotent. Idempotent transformations are discussed in~\cite[Sec.~2.7]{GanyushkinMazorchuk2008} 
where \crefitem{prelim:prop:idempotent}{prop:idemp:item:img} is stated
and it is shown that $\mathcal T_n$ contains 
\[
\sum_{k=1}^n \binom{n}{k} k^{n - k}
\]
idempotent transformations. In~\cite{Beaudry1989CharacterizationIdempotentTM} 
algorithmic and complexity theoretical questions of deciding
whether all transformations of a given transformation semigroup
are idempotent were investigated.

A classical result by 
Howie~\cite{Howie66} states that the set of all simple idempotents on a finite set
generates the transformation semigroup consisting of all non-permutations.
Gomes and Howie~\cite{GomesHowie1987} 
proved that $\binom{n}{2}$ idempotents are sufficient and necessary
to generate this set.